\unnumberedchapter{Product Discovery}\label{ch:product-discovery}

\chaptersubtitle{Deciding what is worth building --- before you build it}

\chapterrule

The most expensive mistake in software is not bad code. It is building the wrong thing \emph{well}~---
pouring months of careful engineering into a product, feature, or system that nobody actually
needs. Slow or buggy code can be fixed; a beautifully built solution to a problem no one has is
simply waste.

\textbf{Product discovery} is the work that happens \emph{before} planning and code: figuring out what is
worth building, for whom, and why~--- and validating that cheaply, before committing to expensive
engineering. It is the most commonly skipped phase of building software, and the one that decides
whether everything downstream even matters. You can plan, code, and operate flawlessly and still
fail completely if discovery was skipped.

This part is about answering one question well: \textbf{should we build this at all~--- and if so, what
exactly?}

\unnumberedsection{true}{Why Discovery Exists: The Cost of Building the Wrong Thing}\label{product-discovery:why-discovery-exists-the-cost-of-building-the-wrong-thing}

Engineering effort is expensive and largely irreversible: once a team has spent three months
building something, the sunk cost makes it hard to admit it was the wrong thing. Discovery exists
because the cost of \emph{finding out you were wrong} should be paid in days of conversation and cheap
experiments, not months of code.

Two ideas frame the whole discipline:

\begin{itemize}
\tightlist
\item
  \textbf{Outputs are not outcomes.} An \emph{output} is something you ship (a feature, a screen, an API). An
  \emph{outcome} is a change in the world (a user succeeds at something, a metric moves, a problem goes
  away). Shipping is easy; creating an outcome is the actual job. A team can ship constantly and
  create no value at all.
\item
  \textbf{The build trap.} Most teams measure themselves by how much they ship, and slide into being a
  ``feature factory''~--- busy, productive-looking, and disconnected from whether any of it matters.
  Discovery is the discipline that keeps effort tied to outcomes instead of motion.
\end{itemize}

\noindent{}The mindset shift discovery asks for: fall in love with the \textbf{problem}, not your solution. Your
first idea is a guess. Discovery is how you find out~--- cheaply~--- whether it's a good one.

\unnumberedsection{true}{The Four Risks Every Idea Must Survive}\label{product-discovery:the-four-risks-every-idea-must-survive}

A useful way to think about discovery (from Marty Cagan) is that every idea carries four risks, and
discovery is the work of \emph{retiring} them before you build:

\begin{itemize}
\tightlist
\item
  \textbf{Value risk}~--- will anyone actually want this? Will they use it, switch to it, pay for it? This
  is the biggest risk and the most often ignored, because it's the easiest to assume away.
\item
  \textbf{Usability risk}~--- can users figure out how to use it? A valuable thing nobody can operate
  fails just as surely as one nobody wants.
\item
  \textbf{Feasibility risk}~--- can \emph{we} actually build it, with our time, technology, data, and skills?
  (This is the one engineers naturally focus on~--- and often the \emph{least} risky of the four.)
\item
  \textbf{Business viability risk}~--- does it work for \emph{our} business? Legal, financial, brand, sales,
  support, and ethical constraints can all kill an otherwise good idea.
\end{itemize}

\noindent{}The trap is to spend all your energy on feasibility (``can we build it?'') while assuming value
(``they'll want it''). Discovery deliberately attacks the \emph{value} risk first, because it's both the
largest and the cheapest to test before any code exists.

\unnumberedsection{true}{Understanding the Problem and the User}\label{product-discovery:understanding-the-problem-and-the-user}

You cannot validate a solution to a problem you don't understand. Discovery starts with the people
you're building for.

\subsection*{Talk to users --- but listen, don't pitch}\phantomsection\label{product-discovery:talk-to-users--but-listen-dont-pitch}

The single highest-leverage discovery activity is talking to real users~--- and the single most
common way to do it badly is to pitch your idea and collect compliments. People are polite; they
will tell you your idea is ``interesting,'' and you will learn nothing.

The discipline (captured well in Rob Fitzpatrick's \emph{The Mom Test}): \textbf{ask about their life and their past behavior,
not your idea.} Good questions are concrete and backward-looking~--- ``Walk me through the last time
you ran into this''~--- because what people \emph{actually did} is data, while what they \emph{say they would do}
is a guess. Avoid leading questions, avoid hypotheticals, and treat compliments as a warning that
you've started pitching.

\subsection*{Find the problem beneath the solution}\phantomsection\label{product-discovery:find-the-problem-beneath-the-solution}

Users almost always arrive with a \emph{solution}, not a problem: ``I need a faster horse,'' ``just add an
export button.'' Your job is to dig past the requested solution to the underlying need~--- the ``five
whys'' is a blunt but effective tool here. The export button might really be ``I need to share this
with my boss, who lives in spreadsheets,'' which a dozen better solutions could serve. Solving the
stated solution literally is how you build features no one uses.

\subsection*{Jobs To Be Done}\phantomsection\label{product-discovery:jobs-to-be-done}

A sharp lens for ``what is the user really trying to do'' is \textbf{Jobs To Be Done}: people ``hire'' a
product to make progress in a particular situation. They don't buy a drill because they want a
drill; they hire it to get a hole~--- and really, to hang a picture. Focusing on the \emph{job} (the
progress someone is trying to make) rather than the demographic or the feature keeps you anchored to
the outcome the user cares about.

\unnumberedsection{true}{Validating That It's Worth Solving}\label{product-discovery:validating-that-its-worth-solving}

Understanding a problem isn't enough~--- plenty of real problems aren't worth solving. A problem is
worth building for when it is \textbf{real, frequent, painful, and something people will change behavior
(or pay) to fix.} A rare, mild annoyance is not a business.

Two practical tools:

\begin{itemize}
\tightlist
\item
  \textbf{The Riskiest Assumption Test.} List the assumptions your idea depends on, find the one that~---
  if false~--- kills the whole thing, and test \emph{that} first, as cheaply as possible. Usually it's a
  value assumption (``people will pay for this''), not a feasibility one. Spending your first effort
  anywhere else is motion, not progress.
\item
  \textbf{Demand before build.} You can measure demand before writing the product: a landing page that
  describes the offer and counts sign-ups, a ``fake door'' (a button for a feature that doesn't exist
  yet, measuring who clicks), or a pre-sell. Real behavior~--- a click, an email, a pre-order~--- is far
  stronger evidence than any survey.
\end{itemize}

\unnumberedsection{true}{Shaping and Testing the Solution Cheaply}\label{product-discovery:shaping-and-testing-the-solution-cheaply}

Once the problem is validated, discovery shifts to the solution~--- and the goal is to learn the most
while building the least.

\subsection*{Define success as an outcome, not an output}\phantomsection\label{product-discovery:define-success-as-an-outcome-not-an-output}

Before designing anything, state what success \emph{is}~--- and state it as an outcome. ``Ship a
recommendations feature'' is an output and tells you nothing about whether it worked. ``Increase the
share of users who find a relevant item in their first session'' is an outcome you can actually test
a solution against. Outcome-framing is what keeps the rest of discovery honest.

\subsection*{Prototype and test before building}\phantomsection\label{product-discovery:prototype-and-test-before-building}

You can test a solution long before you build it, with experiments that cost hours instead of
months:

\begin{itemize}
\tightlist
\item
  \textbf{Prototypes}~--- paper sketches or clickable mockups put in front of real users to test usability
  and desirability.
\item
  \textbf{Wizard of Oz}~--- the interface looks real, but a human does the work behind the curtain, so you
  test the experience without building the engine.
\item
  \textbf{Concierge}~--- you deliver the outcome manually for a few real users (no product at all) to learn
  whether it's wanted and what it really takes.
\item
  \textbf{Fake door / landing page}~--- measure interest in a thing that doesn't exist yet.
\end{itemize}

\noindent{}Each retires a specific risk (usability, value, feasibility) for a fraction of the cost of building.

\subsection*{The MVP, correctly understood}\phantomsection\label{product-discovery:the-mvp-correctly-understood}

The \textbf{Minimum Viable Product} is the most misunderstood term in software. It is \textbf{not} a
low-quality first version of your product. It is the \emph{smallest thing that lets you test your core
hypothesis with real users}~--- an experiment, not a deliverable. The ``viable'' means \emph{viable as a
test of the idea}, not \emph{viable as a product to live with}. If you can learn what you need from a
landing page or a concierge run, that is your MVP, and you've saved months.

\unnumberedsection{true}{Prioritization: Deciding What to Pursue}\label{product-discovery:prioritization-deciding-what-to-pursue}

Discovery surfaces more opportunities than you can pursue, so you need a way to choose. A few
frameworks, used as thinking aids rather than oracles:

\begin{itemize}
\tightlist
\item
  \textbf{Value vs effort}~--- the simplest: plot ideas by expected value against cost, and start in the
  high-value / low-effort corner. Crude, but it filters out a lot.
\item
  \textbf{RICE}~--- score by \textbf{R}each × \textbf{I}mpact × \textbf{C}onfidence ÷ \textbf{E}ffort, to compare opportunities
  on a consistent (if rough) basis.
\item
  \textbf{MoSCoW}~--- sort into Must-have / Should-have / Could-have / Won't-have to force explicit
  trade-offs and an explicit ``not now.''
\item
  \textbf{Kano model}~--- distinguishes \emph{basic} expectations (their absence angers, their presence is
  unnoticed), \emph{performance} features (more is better), and \emph{delighters} (unexpected joys). Useful for
  deciding where effort actually changes how users feel.
\item
  \textbf{Opportunity Solution Trees} (Teresa Torres, \emph{Continuous Discovery Habits})~--- map a desired \textbf{outcome}
  to the \textbf{opportunities} (user needs / pain points) that could move it, then to candidate
  \textbf{solutions}, then to \textbf{experiments}. It keeps a whole team's ideas tied back to the outcome
  instead of to a feature wish-list.
\end{itemize}

\noindent{}The point of any of these is not the number they produce; it's that they force the conversation and
make the reasoning behind a choice explicit.

\unnumberedsection{true}{A Worked Example: Discovery for a Notes Service}\label{product-discovery:a-worked-example-discovery-for-a-notes-service}

Suppose a team wants to build a small notes service~--- the same Notes API that Part IV of this book
takes all the way to production. Discovery for it might look like this:

\begin{itemize}
\tightlist
\item
  \textbf{Problem, not solution.} Interviews with a dozen developers reveal the real pain is not ``I need
  a notes app'' but ``the snippets and decisions I jot down during the day are scattered across chat,
  sticky notes, and scratch files, and I cannot find them a week later.''
\item
  \textbf{Riskiest assumption.} \emph{Value risk}: will people actually send notes to an API from their own
  tools, or will they keep using whatever is already open on their screen?
\item
  \textbf{Smallest test.} Before writing a server, give five developers a two-line shell alias that
  appends a note to a shared file, and watch for two weeks: do they use it daily, and do they
  come back to search it?
\item
  \textbf{Outcome.} Success is ``developers retrieve a note they saved at least once a week,'' not ``we
  shipped create, read, update, and delete endpoints.''
\item
  \textbf{Decision.} If usage stalls after the first days, stop~--- the evidence says the problem is not
  painful enough. If it holds, the validated problem, the target users, and the success measure
  become the input to project planning.
\end{itemize}

\noindent{}Notice that nothing here required the API to exist. Discovery retires the biggest risk with the
cheapest possible experiment, and only then does building begin.

\unnumberedsection{true}{Common Traps}\label{product-discovery:common-traps}

Discovery fails in predictable ways. Watch for:

\begin{itemize}
\tightlist
\item
  \textbf{Solution-first thinking}~--- falling in love with an answer and reverse-justifying the problem.
\item
  \textbf{The feature factory}~--- measuring success by shipping volume, never by outcomes.
\item
  \textbf{Confirmation bias}~--- running ``tests'' designed to hear yes; talking only to people who already
  agree; counting compliments as validation.
\item
  \textbf{Building because you can}~--- technical excitement (``we could use this new model / framework'')
  masquerading as user value.
\item
  \textbf{Skipping discovery under deadline pressure}~--- the most expensive shortcut there is, because it
  trades a few days of cheap learning for months of possibly wasted building.
\item
  \textbf{Mistaking activity for progress}~--- interviews, decks, and experiments that never actually
  retire a real risk.
\end{itemize}

\noindent{}The antidote to all of them is the same: tie every step to a real, falsifiable question about value,
and be willing to hear ``no.''

\unnumberedsection{true}{The Handoff to Project Planning}\label{product-discovery:the-handoff-to-project-planning}

Discovery is done when you can answer, with evidence rather than hope: \emph{who} this is for, \emph{what}
problem it solves, \emph{why} it's worth solving, what \emph{outcome} defines success, and which \emph{riskiest
assumptions} you've retired. At that point the question flips from ``should we build this?'' to ``what
exactly are we building, and how will we document and plan it?''

That is precisely where \textbf{\hyperref[ch:project-planning-guide]{Project Planning}}
(Part II) picks up: it turns a validated, well-understood problem into a problem statement,
architecture, and the set of documents a real project keeps. Discovery decides \emph{whether and what};
planning decides \emph{exactly what and how}. Skip discovery and you plan beautifully for the wrong
thing; skip planning and you build the right thing chaotically.

\unnumberedsection{true}{A Discovery Checklist}\label{product-discovery:a-discovery-checklist}

Before committing engineering effort, you should be able to answer:

\begin{itemize}
\tightlist
\item
  \textbf{Problem}~--- what real, frequent, painful problem am I solving, and for whom? (Stated as a
  problem, not a solution.)
\item
  \textbf{Evidence}~--- have I talked to real users about their actual behavior, not pitched my idea?
\item
  \textbf{Value}~--- is there evidence people will \emph{change behavior} (use, switch, pay) to get this?
\item
  \textbf{Outcome}~--- what observable outcome defines success, and how will I measure it?
\item
  \textbf{Riskiest assumption}~--- what single assumption, if wrong, kills this~--- and have I tested it
  cheaply?
\item
  \textbf{Smallest test}~--- what is the least I can build or fake to learn whether this works?
\item
  \textbf{Decision}~--- am I prepared to \emph{not build it} if the evidence says no?
\end{itemize}

\noindent{}If you can answer these honestly, you've done the work that makes everything downstream worth doing.
